\Qhead{3.9}{Рассеяние волн в случайно-неоднородных средах}{Рассеяние волн в случайно-неоднородных средах. Борновское приближение, метод плавных возмущений. Рассеяние волн на шероховатой поверхности. Понятие об обратной задаче рассеяния.}
\Subfile{3.9.1}{sections/03_statistics/q09/s01}
\Subfile{3.9.2}{sections/03_statistics/q09/s02}

\recap{Слабое рассеяние связывает пространственный спектр неоднородностей с угловым спектром поля. Борн разлагает само поле, плавные возмущения — его комплексный логарифм. Обратная задача ограничена полосой, апертурой и потерей фазы.}
\board{Интегральное уравнение, первая итерация, переданный вектор $\bm q$; фазовый набег шероховатой границы.}
\oral{Почему уменьшение зеркально отражённой мощности ещё не доказывает поглощение?}
\bait{Можно оставить потерю фазы: одинаковые интенсивности не обязаны однозначно задавать комплексную амплитуду рассеяния и объект.}
\sources{Barbastathis~\cite{barbastathis}, введение в обратные задачи; Сорокин~\cite{sorokinray}, отражение от поверхности. Формулы слабого объёмного рассеяния и гауссовской шероховатости получены здесь из интегрального уравнения и среднего фазового множителя; рекомендованная программа для углубления — Рытов, Кравцов, Татарский, ч.~2 (полный текст этого издания не проверен).}
